\section{总结与延伸}

\subsection{全课知识图谱}

本课建立了从传统方法到神经网络方法的完整语言建模认知链：

\begin{center}
\begin{tikzpicture}[node distance=1.0cm, every node/.style={draw, rounded corners, minimum width=4cm, minimum height=0.8cm, font=\small}]
\node (basics) {深度学习实用技巧};
\node (lm) [below=of basics] {语言模型概念};
\node (ngram) [below=of lm] {N-gram 语言模型};
\node (fwnn) [below=of ngram] {固定窗口神经语言模型};
\node (rnn) [below=of fwnn] {循环神经网络 (RNN)};
\node (grad) [below=of rnn] {梯度消失与爆炸};
\draw[-{Stealth}, thick] (basics) -- (lm);
\draw[-{Stealth}, thick] (lm) -- (ngram);
\draw[-{Stealth}, thick] (ngram) -- (fwnn);
\draw[-{Stealth}, thick] (fwnn) -- (rnn);
\draw[-{Stealth}, thick] (rnn) -- (grad);
\node[draw=none, right=0.8cm of basics, text width=5.5cm, font=\scriptsize] {Dropout、Xavier 初始化、Adam};
\node[draw=none, right=0.8cm of lm, text width=5.5cm, font=\scriptsize] {$P(w \mid \text{context})$、链式法则};
\node[draw=none, right=0.8cm of ngram, text width=5.5cm, font=\scriptsize] {计数、马尔可夫假设、稀疏性};
\node[draw=none, right=0.8cm of fwnn, text width=5.5cm, font=\scriptsize] {词嵌入拼接、无稀疏性、仍受限};
\node[draw=none, right=0.8cm of rnn, text width=5.5cm, font=\scriptsize] {参数共享、隐藏状态、BPTT};
\node[draw=none, right=0.8cm of grad, text width=5.5cm, font=\scriptsize] {梯度裁剪、LSTM/GRU $\rightarrow$ Transformer};
\end{tikzpicture}
\end{center}

\subsection{关键 Takeaways}

\begin{importantbox}{五条核心要点}
\begin{enumerate}
    \item \textbf{语言模型是 NLP 的基石}：从手机输入法到 ChatGPT，语言模型的核心任务始终是预测下一个词的概率分布
    \item \textbf{从计数到学习}：N-gram 靠统计计数，神经语言模型靠参数学习——后者在泛化能力上有质的飞跃
    \item \textbf{RNN 的核心创新是参数共享}：同一组权重在每个时间步重复应用，使模型能处理任意长度的输入
    \item \textbf{梯度问题是序列模型的命门}：梯度消失限制了 RNN 学习长距离依赖的能力，这一问题推动了 LSTM、GRU 乃至 Transformer 的诞生
    \item \textbf{实用技巧至关重要}：Dropout、合理初始化、Adam 优化器、梯度裁剪——这些``小技巧''的组合是深度学习成功的关键
\end{enumerate}
\end{importantbox}

\subsection{展望：从 RNN 到 Transformer}

Manning 在课程最后暗示了后续的发展方向：

\begin{itemize}
    \item \textbf{LSTM 和 GRU}：通过门控机制缓解梯度消失问题，使 RNN 能更好地捕捉长距离依赖（下一讲）
    \item \textbf{编码器-解码器架构}：将 RNN 组合为翻译模型、摘要模型等序列到序列系统
    \item \textbf{注意力机制}（Attention）：让解码器在每步直接``看到''编码器的所有隐藏状态，而非仅依赖最终状态
    \item \textbf{Transformer}：完全抛弃循环结构，用自注意力实现全并行计算，彻底解决 RNN 的顺序计算瓶颈和长距离依赖问题
\end{itemize}

\subsection{拓展阅读}

\begin{itemize}
    \item Bengio et al., ``A Neural Probabilistic Language Model'' (2003): \url{https://www.jmlr.org/papers/volume3/bengio03a/bengio03a.pdf} —— 首个神经语言模型
    \item Mikolov et al., ``Recurrent Neural Network based Language Model'' (2010) —— RNN 语言模型的经典论文
    \item Hochreiter \& Schmidhuber, ``Long Short-Term Memory'' (1997) —— LSTM 的原始论文
    \item Cho et al., ``Learning Phrase Representations using RNN Encoder-Decoder'' (2014) —— GRU 的提出
    \item Karpathy, ``The Unreasonable Effectiveness of Recurrent Neural Networks'' (2015): \url{http://karpathy.github.io/2015/05/21/rnn-effectiveness/} —— RNN 生成文本的有趣博客
    \item CS224N 课程主页：\url{https://web.stanford.edu/class/cs224n/}
\end{itemize}

\end{document}
